\documentclass{report}
\usepackage{amsmath,amssymb}
\setlength{\parindent}{0mm}
\setlength{\parskip}{1em}
\begin{document}
\begin{center}
\rule{6in}{1pt} \
{\large Eric de Sturler \\
{\bf Randomized Methods and Model Reduction for Accelerating the Solution of Inverse Problems}}

Department of Mathematics \& \\ Computational Modeling and Data Analytics \\ Academy of Integrated Sciences \\ 460 McBryde \\ Virginia Tech \\ 225 Stanger Street \\ Blacksburg \\ VA 24061-0123 \\ USA
\\
{\tt sturler@vt.edu}\\
Selin Sariaydin \\
Misha Kilmer\\
Drayton Munster\end{center}

In nonlinear inverse problems, the objective function often involves the
solution of a discretized PDE for many right hand sides, corresponding to
many measurements.
Additional linear systems must be solved for evaluating or approximating
the Jacobian of the nonlinear least squares problem. Hence, the solution
of the inverse problem
requires the solution of a very large number of large linear systems. We
discuss randomized techniques as well as model reduction approaches to
drastically reduce the number of systems to be solved. We apply our
approach to problems in diffuse optical tomography.


\end{document}
